\documentclass[sigconf, anonymous]{acmart}

%% Using the option "anonymous" will produce "Anonymous Author(s)"
%% for double-anonymized review. Upon acceptance, remove this option
%% and set the authors' names and affiliations using the corresponding
%% commands.

%% CBSoft uses an ACM-like format without the ACM Reference block.
%% The following warning from acmart can be safely ignored:
%% "ACM reference format is mandatory for papers over one page."
\settopmatter{printacmref=false}
\renewcommand\footnotetextcopyrightpermission[1]{} % removes footnote with conference information in first column

%% CBSoft uses an ACM-like format without a copyright note.
\setcopyright{none}

%% These commands are for a PROCEEDINGS abstract or paper.
\acmConference[SBCARS 2026]{20th Brazilian Symposium on Software Components, Architectures, and Reuse}{September 8--11, 2026}{São Paulo, SP, Brazil}

% CBSoft uses an ACM-like format without CCS Concepts.
% The following warning from acmart can be safely ignored:
% "CCS concepts are mandatory for papers over two pages."

%% \BibTeX command to typeset BibTeX logo in the docs
\AtBeginDocument{
    \providecommand\BibTeX{{Bib\TeX}}
}

%% Only for papers written in PORTUGUESE: 
%% comment out the following three lines
%\renewcommand\abstractname{Resumo}
%\renewcommand\keywordsname{Palavras-chave}
%\renewcommand\refname{Referências}

%% end of the preamble, start of the body of the document source.


% ------------------------------------------------------------
% Authors: please make sure to read the checklist comments
% provided at the end of this document before submission.
% ------------------------------------------------------------
\begin{document}

%% The "title" command has an optional parameter,
%% allowing the author to define a "short title" to be used on the page 
%% headers.
\title{A Rapid Review on Observability of Cloud-Native Applications using Kubernetes}

%% The "author" command and its associated commands are used to define
%% the authors and their affiliations.
%% Of note is the shared affiliation of the first two authors, and the
%% "authornote" and "authornotemark" commands
%% used to denote shared contribution to the research.
\author{Fagner Silva}
\affiliation{
  \institution{Centro de Informática, UFPE}
  \city{Recife}
  \country{Brasil}
}
\email{ffcs@cin.ufpe.br}

\author{Jackson Raniel}
\affiliation{
  \institution{Universidade de Pernambuco - Campus Caruaru}
  \city{Caruaru}
  \country{Brasil}
}
\email{jackson.florencio@upe.br}

\author{Vinícius Cardoso Garcia}
\affiliation{
  \institution{Centro de Informática, UFPE}
  \city{Recife}
  \country{Brasil}
}
\email{vcg@cin.ufpe.br}

%% By default, the full list of authors will be used on the page
%% headers. This list is often too long and will overlap
%% other information printed in the page headers. 
%% This command allows the author to define a more concise list
%% of authors' names for this purpose.
\renewcommand{\shortauthors}{Surname1stAuthor et al.}
\renewcommand{\shorttitle}{A Rapid Review on Observability of CNA using K8S}

%% The abstract is a short summary of the work the paper presents.
\begin{abstract}
The rise of cloud-native applications and the widespread adoption of Kubernetes have introduced significant complexity in system management. In these dynamic environments, traditional monitoring is insufficient, making observability essential for understanding internal system states through metrics, logs, traces, and events. This paper analyzes the state of the art in observability for Kubernetes-based environments, focusing on open-source tools and architectural patterns to minimize vendor lock-in and costs. We employed a Rapid Review methodology, following a structured protocol to synthesize evidence from major scientific databases published between 2018 and October 2025. The study screened relevant literature to identify tools, challenges, and architectural gaps. The findings highlight the lack of a widely adopted formal reference architecture capable of holistically managing the observability data lifecycle. Instead, a '\textit{de facto}' pattern has emerged, integrating specific open-source tools—primarily Prometheus, Loki, Jaeger, and Grafana—with OpenTelemetry becoming the standard for unified telemetry representation. However, the fragmentation of these tools creates significant operational complexity and a high cognitive load for engineering teams. We conclude that, while a mature open-source ecosystem exists, it lacks a unified platform, requiring manual integration. Future directions include adopting AIOps for automated anomaly detection and integrating security (SecOps) to manage the increasing volume and multimodal nature of telemetry data.
\end{abstract}

%% Keywords. The author(s) should pick words that accurately describe
%% the presented work. Separate the keywords with commas.
\keywords{Observability, Metrics, Logs, Tracing, Kubernetes, Open-source, Cloud-native.}

%% This command processes the author, affiliation, and title
%% information and builds the first part of the formatted document.
%% "frenchspacing" avoids an additional space after a period at the end of a sentence.
\frenchspacing
\maketitle

\section{Introduction}

The rise of cloud-native applications, built to leverage the distributed capabilities of the cloud delivery model, is a defining feature of modern software engineering \cite{vancamp24, oracle_2021}. Solidified by paradigms like the 12-Factor App\footnote{\textit{The Twelve-Factor App} is a methodology for building modern SaaS applications. Available at: \url{https://12factor.net/}.} and orchestrated almost universally by Kubernetes (K8s), this architectural style enables applications to operate at a global scale with unprecedented resilience \cite{gannon_2017}. However, this distribution and dynamism come at a cost: complexity. Unlike monolithic systems, cloud-native applications are composed of many ephemeral, containerized microservices, making fault detection and performance analysis ``orders of magnitude more complex'' \cite{usman_2022}.

This inherent complexity renders traditional monitoring approaches insufficient. In K8S-orchestrated environments, ``rapid scalability, the ephemerality of containers, and service interdependencies can obscure the root causes of problems.'' Consequently, \textbf{observability}---the practice of understanding a system's internal state from its external outputs \cite{kalman_1959}---has emerged as an essential discipline for managing these systems. It relies on the comprehensive collection and analysis of metrics, logs, and distributed traces to enable effective troubleshooting and performance optimization.

In response to this need, a powerful \textit{de facto} standard based on open-source tools, largely from the Cloud Native Computing Foundation (CNCF), has been widely adopted. However, while a formal, unified observability architecture for cloud-native systems remains elusive in academic literature \cite{araujo_24, zanatta_23, kosinska_2023}, the industry has coalesced around this community-driven stack. The problem, which this paper addresses, is that this \textit{de facto} solution is not a single, cohesive platform but rather a fragmented ecosystem of independent tools.

This paper argues that this fragmentation, while offering flexibility, introduces a significant, often underestimated, operational burden and architectural trade-offs. The core contribution of this work is a critical analysis of the consequences of this \textit{de facto} standard. Through a Rapid Review of current practices and literature, we: 1) dissect the components of this common stack, 2) identify the integration gaps and operational challenges it creates, and 3) evaluate the architectural implications of relying on a loosely coupled toolchain for a mission-critical function. This analysis provides a clear perspective on the limitations of the current state-of-the-art and sets the stage for future research into more integrated and intelligent observability platforms. This study is structured as follows: background and related work, methodology, findings and results, and conclusion.

\section{Background and Related Work}
\label{sec:background}

The definition of the cloud-native term is usually related to something built on the cloud, rather than being on-premise \cite{gannon_2017}. However, this definition is incomplete, as cloud-native refers to the software method of creating, deploying, and managing modern applications in cloud computing environments \cite{aws_2025}. Therefore, it is possible to host cloud-native applications (CNAs) both in the cloud and on-premises, where both styles are sometimes used simultaneously.

The proliferation of cloud-native applications, in conjunction with this hosting model, has introduced a new level of complexity in analyzing problems within the underlying infrastructure. Therefore, as a direct consequence, the complexity of tracking application performance, identifying bottlenecks, and ensuring system reliability has increased exponentially \cite{sakinala_2025}. Where organizations once relied on centralized, monolithic systems, modern environments feature distributed frameworks with containerized workloads, independently deployable microservices, and multi-cloud infrastructure.

On the other hand, the term observability is not a concept that originated early in the field of software engineering \cite{kalman_1959}. However, several related terms are used interchangeably when referencing this concept \cite{kosinska_2023}. According to \cite{songtab_1984}, observability is a fundamental system-theoretic property that reflects the possibility of estimating internal states on the basis of input/output data. Therefore, it is possible to determine the system's operational viability through these external indicators.

According to \cite{kosinska_2023}, observability is often defined as consisting of three pillars: (i) metrics (collected while monitoring activities), (ii) logging, and (iii) tracing. However, as mentioned by \cite{hartikainen_2024}, the mere fact that a system possesses these three pillars does not automatically make it observable. The most value is obtained by combining these pillars and converting them into interpretable data.

Later, \cite{bhatia_2025} informs that an additional pillar is used in modern observability frameworks: events. These essential telemetry categories, when combined, facilitate comprehensive system visibility. For further clarification, the definitions of each concept will be provided in the next section.

Metrics represent the raw measurements of resource usage or behavior that can be observed and collected throughout your systems \cite{ellingwood_2017}. They can be interpreted as counters or measurements of a system characteristic during a time period; they are numeric by definition and represent aggregated data. Examples of metrics include the number of requests per second responded to by the HTTP server and the total heap memory allocated to the Java microservice.

Logs are the primary data source for developers and software engineers to uncover the underlying problems in computing systems. Logs are capable of exporting information about each operation and the system state, making it possible to understand exactly what a process was doing at a given point in time \cite{sre_2016}.

A distributed trace, more commonly known as a trace, records the paths taken by requests (made by an application or end-user) as they propagate through multi-service architectures, like microservice and serverless applications \cite{otel_primer_2025}. A trace is composed of one or more spans. A span represents a unit of work or operation. Spans track specific operations that a request performs, providing a clear picture of what happened during the time the operation was executed.

Events can be utilized to enhance a system's observability \cite{hartikainen_2024}. They are registered with specific actions, such as the execution of a function, the updating of a database record, or an exception thrown by the code. They are timestamped, unchangeable documentation that comes in three forms: plain text, structured, and binary.

After understanding the core tenets of observability, it is necessary to make a clear distinction between monitoring and observability.

Although often used interchangeably, monitoring and observability are distinct concepts. Monitoring tracks a system's health using predefined metrics and logs to detect known failures. However, this approach can overlook unforeseen issues in dynamic, distributed systems, as they do not fit into predefined problem categories. Observability, in contrast, is the ability to investigate and determine where and why a failure occurred, especially for these unknown issues.

Therefore, monitoring is a prerequisite for observability, as the two are closely intertwined. A simple way to view it is that monitoring is the alarm that warns of an existing problem, while observability provides the tools to discover why the problem occurred. The former alerts operators to a failure, and the latter helps identify its root cause, as shown in Figure \ref{fig:obs-vs-mon}.

\begin{figure}[ht]
    \centering
    \includegraphics[width=0.75\linewidth]{images/monitoring-vs-observability.png}
    \caption{Observability vs Monitoring}
    \label{fig:obs-vs-mon}
\end{figure}

There is a popular saying in software engineering: if something cannot be measured, it cannot be controlled. Whenever a system needs to be monitored, an instrument or an instrumentation system must be in place to facilitate this process. For clarification, an instrument (or instrumentation system) is a device that measures a specific quantity \cite{ribbens_2013}.

In the purest sense, telemetry is a technology that allows data acquired in one location to be monitored and displayed in another \cite{dondelinger_2013}. The word telemetry is formed by combining two Greek words: \textit{tele}, meaning remote, and \textit{metron}, meaning measure.

In cloud-native environments, telemetry has become more abundant and diverse, presenting both a challenge and an opportunity to go beyond monitoring and achieve true observability, which is necessary to manage complexity. Such complexity is further enhanced through the process of instrumenting CNAs, which can be achieved using different levels of instrumentation. 

Considering the existence of multiple levels of instrumentation, it is possible to implement the instrumentation process across several components, i.e., from the application level to the platform level, where K8S is installed.

The first level is the code/application level. In this level, instrumentation is implemented directly in the application's source code. Developers use language-specific libraries or SDKs (Software Development Kits) (e.g., OpenTelemetry SDK, Micrometer, Prometheus Client Libraries) to define and expose metrics, create structured logs, and initiate/propagate tracing spans.

The second level is the platform level. This level abstracts network instrumentation away from the application. A sidecar proxy (e.g., Envoy, Linkerd) is injected into each application pod. This proxy intercepts all network traffic (incoming and outgoing) from the application container. Service Mesh (e.g., Istio) manages and configures these proxies to collect telemetry in a standardized way.

The third level is the host/node level. A software agent runs as a DaemonSet (one instance per cluster node). This agent collects metrics and logs from the node's operating system and running containers, typically reading information from virtual file systems such as Prometheus Node Exporter and cAdvisor.

The final level is the kernel level. The extended version of the Berkeley Packet Filter (eBPF) is a Linux kernel technology that allows programs to be run safely and efficiently in a sandbox within the kernel itself. Without altering application or kernel code, these programs can hook into low-level instrumentation points, such as system calls, network functions, and process scheduler events. Therefore, it can be classified as an instrumentation at the kernel level in Linux \cite{rise_2023}.

\subsection{Related Work}

Regarding related work, we highlight the work of \cite{kosinska_2023}, which seeks to chart a course for the so-called state of the art in observability for cloud-native applications through a systematic literature review spanning 2018 to 2022. That work aims to detail various observability approaches, classifying their application by management area (DevOps and testing), requirement type (functional and non-functional), and ecosystem (infrastructure or microservices). Our work focuses on more pressing matters related to achieving a minimum level of observability using open-source tools in a cloud-native environment.

On the other hand, the work of \cite{usman_2022} is more closely aligned with our work, although it does not propose a solution to the aforementioned issues. Its focus is divided into two parts: the first seeks to map aspects of the state of the art in observability for distributed environments and microservices. The second part explores which components of data telemetry can be used to achieve an ideal state of observability. However, there is a lack of concern around practical aspects, such as designing and evaluating an observability framework. This study leverages these core aspects of observability, and it reflects on different ways to achieve a viable framework.

Another line of study \cite{faseeha_2025} classifies existing solutions according to architectural patterns such as gateway, service mesh, and event-driven. Its focus emphasizes how observability research is evolving through the use of new frameworks. However, no practical approach was proposed that considered the observability pillars. Its primary concern is microservice architecture, and there is confusion between what is called `monitoring parameters' and those pillars.

\section{Methodology}
\label{sec:methodology}

Rapid Reviews (RRs) are secondary studies designed to provide timely research evidence to support decision-making in practical environments \cite{cartaxo_2018}. RRs are conducted considering constraints inherent to practical environments, such as time and effort, thus delivering evidence more promptly and at lower costs, often communicated through accessible mediums.

This approach is particularly suitable for our study, as it aims to swiftly synthesize scientific evidence on the rapidly evolving field of observability in cloud-native applications without compromising systematic rigor. An RR complements Systematic Reviews (SRs) by filling an essential gap in promoting evidence-based, practice-oriented software engineering for urgent research questions.

Similar to SRs, this Rapid Review consisted of three main phases: planning, performing, and reporting. These phases are detailed in the following subsections.

\subsection{Planning Phase}
In the planning phase, we followed a detailed protocol for this Rapid Review, as established by Cartaxo \textit{et al.} \cite{cartaxo_2018}. This protocol encompassed the formulation of the research problem, the specific research questions (RQs), the search strategy, the selection criteria, and the data extraction procedure. The online tool Parsifal was utilized to systematically compile and manage all aspects of this protocol, ensuring transparency and adherence to a predefined methodology.

Cloud-native environments, characterized by microservices, containers, and dynamic orchestration (e.g., Kubernetes), inherently generate an exponentially growing volume, variety, and velocity of telemetry data (logs, metrics, traces, and events). Managing, processing, and storing this vast amount of raw data becomes economically prohibitive and technically unfeasible. As highlighted by \cite{sridharan_2018}, this exponential growth in data significantly constrains a system's capacity to process information and to deal effectively with ``unknown unknowns,'' necessitating robust observability solutions.

\subsection{Research Questions}
The goal of this study is to identify the state of the art in observability of Kubernetes-based cloud-native applications and investigate the open-source tools and architectural patterns deployed to achieve the degree of observability proposed in the literature. To guide this Rapid Review, the following research questions were formulated:

\begin{itemize}
    \item \textbf{RQ01:} What are the essential components of an observability architecture for Kubernetes-based cloud-native applications?
    \item \textbf{RQ02:} What reference architectures or observability frameworks for Kubernetes-based cloud-native environments have been proposed in the literature?
    \item \textbf{RQ03:} What open-source tools are used or recommended to implement the different components of a Kubernetes-based observability architecture for cloud-native applications?
    \item \textbf{RQ04:} What research gaps and future directions have been identified in the literature on open-source observability architectures for Kubernetes-based cloud-native applications?
\end{itemize}

\subsection{Performing Phase}
The performing phase involved the systematic execution of the search and study selection procedures as defined in the planning phase. The literature search was primarily conducted using four leading scientific databases known for their extensive coverage in Computer Science and Software Engineering: IEEE Xplore, ACM Digital Library, SpringerLink, and ScienceDirect. These databases were selected for their relevance to cloud computing, distributed systems, and software observability. The temporal scope of the search was set from January 2018 to October 2025, with a focus on recent developments in the field.

The search string, carefully constructed and adapted for each database's specific syntax using the filter full-text, was:

\begin{quote}
\textbf{Search String:} (``cloud native" OR ``cloud-native'' OR ``microservices") AND (``observability'' OR ``monitoring'' OR ``tracing'' OR ``logging'' OR ``metrics'' OR ``events'') AND (``kubernetes" OR ``k8s'') AND (``open source" OR ``open-source'')
\end{quote}

Initial searches identified seminal works such as \cite{kosinska_2023} and \cite{usman_2022}, which served as starting points for a complementary forward and backward citation analysis to identify additional relevant studies. Furthermore, an AI-powered literature discovery tool (Research Rabbit) was used to expand the literature search through intelligent recommendation algorithms, efficiently identifying related works and emerging topics based on the initial relevant papers.

The identified studies underwent a multi-stage selection process based on predefined Inclusion and Exclusion Criteria. The study selection process consists of three distinct stages: initial screening, full-text review, and conflict resolution. Initially, titles and abstracts were reviewed by a single researcher to assess their potential relevance. In the second stage, studies that passed the initial screening proceeded to full-text review to confirm their direct relevance to the research questions and ensure full compliance with all criteria. In the third case, when uncertainty or disagreement arose during the full-text review, a second researcher was consulted to resolve conflicts and make a final decision on inclusion.

Six inclusion criteria (IC) were used: \textbf{IC01} - Primary focus on Observability for software systems; \textbf{IC02} - Context related to cloud-native applications or microservices; \textbf{IC03} - Discussion of Kubernetes as the orchestration platform; \textbf{IC04} - Focus on open-source observability tools or architectures; \textbf{IC05} - Articles published from January 2018 to October 2025; \textbf{IC06} - Articles written in English.

Similarly, five exclusion criteria (EC) were used: \textbf{EC01} - Duplicated papers; \textbf{EC02} - Gray literature (e.g., blog posts, company whitepapers); \textbf{EC03} - Focus exclusively on monitoring without addressing observability concepts; \textbf{EC04} - Studies where observability is only tangentially mentioned or not a primary contribution; \textbf{EC05} - Articles exclusively discussing proprietary or closed-source solutions without comparison to open-source alternatives.

For each selected study, data relevant to the research questions were systematically extracted by one researcher and validated by another using the Parsifal online tool. The following information was collected: \textbf{for RQ01:} Definitions of essential architectural components, layers of observability (e.g., data types, functional capabilities, instrumentation levels). \textbf{For RQ02}, descriptions of proposed architectural patterns, frameworks, or significant architectural considerations. \textbf{For RQ03}, the identification of specific open-source tools, their functionalities, and their application within Kubernetes-based environments. Finally, \textbf{for RQ04}, identified research challenges, gaps, and future directions related to open-source observability. This structured approach ensured consistency and scientific rigor in the data collection process.

The reporting phase involved synthesizing and analyzing the extracted data to answer the research questions, identify key findings, and discuss their implications. The subsequent sections of this paper present these findings and outline directions for future work.

\section{Findings and Discussions}
\label{sec:findings}

This section presents a synthesis of the literature review conducted in accordance with the defined research questions. It provides a comprehensive overview of the current state of observability in Kubernetes-based cloud-native applications, identifies prevalent solutions, and highlights research gaps and future directions.

Regarding the essential components of an observability architecture (RQ1), the literature indicates that an effective strategy relies on a combination of core telemetry types and functional capabilities implemented across multiple layers \cite{kosinska_2023, usman_2022, faseeha_2025}. The fundamental pillars are metrics for health overviews, logs for granular event debugging, and distributed traces for tracking request journeys in microservices \cite{gomes_2025, araujo_24}. Beyond these, the literature emphasizes the ``synergistic value'' of cross-analyzing these domains simultaneously to identify ``unknown unknowns'' \cite{usman_2022, araujo_24}. Functional capabilities must span the entire Observability Data Life Cycle (ODLC)—from collection and transmission to storage and visualization \cite{araujo_24, han_2024}. Moreover, modern architectures must provide visibility at the system, service, and network levels \cite{faseeha_2025}, often leveraging low-level instrumentation such as eBPF to ensure a non-intrusive view of both business logic and kernel-level performance \cite{faseeha_2025, usman_2022}.

In terms of proposed architectures and frameworks (RQ2), a significant finding is the scarcity of widely adopted, formal reference architectures specifically tailored for the dynamic nature of Kubernetes. While academic research has proposed structured life cycles and architectural system layers \cite{araujo_24}, there is a noted lack of proposals that adequately manage the simultaneous collection and correlation of metrics, logs, and traces in an integrated manner \cite{araujo_24}. By synthesizing information from studies such as \cite{usman_2022}, it is possible to infer \textit{de facto} architectural patterns driven by the integration of open-source tools. These patterns typically consist of a data collection layer utilizing agents/SDKs, an aggregation layer for centralized processing, a specialized storage layer (e.g., Time Series Databases for metrics), and an analysis layer for dashboarding and root cause analysis \cite{han_2024, araujo_24}.

Regarding open-source tools (RQ3), the review identifies a mature ecosystem where OpenTelemetry (OTel) has emerged as the industry standard for unified instrumentation and telemetry representation \cite{kosinska_2023, faseeha_2025}. Meanwhile, specific tools dominate their respective pillars: Prometheus for metrics, Loki or Fluentd for logging, and Jaeger or Zipkin for distributed tracing, with Grafana serving as the primary visualization hub \cite{usman_2022, faseeha_2025}. Tools like Cilium and Hubble further enhance network-level observability through eBPF-based insights \cite{faseeha_2025}. Despite this maturity, the literature warns that these tools are rarely ``plug-and-play,'' imposing significant operational complexity on engineering teams who must manually manage configurations and data correlation \cite{kosinska_2023, usman_2022}.

Finally, regarding research gaps and future directions (RQ4), the literature highlights the absence of holistic, ``all-in-one'' platforms, forcing practitioners to rely on fragmented stacks \cite{zanatta_23, usman_2022}. A critical challenge remains the correlation of multimodal observability data at scale without losing fine-grained variations \cite{han_2024}. Future research directions emphasize the application of AI and Machine Learning (AIOps) for automated anomaly detection and root cause analysis (RCA) amidst massive data volumes \cite{han_2024, zanatta_23, araujo_24}. Additionally, the integration of security monitoring (SecOps) and the enhancement of Developer Experience (DevEx) are highlighted as essential for reducing the cognitive load associated with managing observability in complex distributed environments \cite{kosinska_2023, deng_2024}.

Table \ref{tab:rq_mapping} compiles all studies mentioned in this paper, along with their relationships to the four research questions discussed.

\begin{table}[ht]
\centering
\caption{Correspondence between Selected Studies and Research Questions}
\label{tab:rq_mapping}
\begin{tabular}{|c|c|c|c|c|}
\hline
\textbf{Ref.} & \textbf{RQ1} & \textbf{RQ2} & \textbf{RQ3} & \textbf{RQ4} \\ \hline
\cite{kosinska_2023}    & \checkmark & \checkmark & \checkmark & \checkmark \\ \hline
\cite{usman_2022}       & \checkmark & \checkmark & \checkmark & \checkmark \\ \hline
\cite{faseeha_2025}     & \checkmark & \checkmark & \checkmark & \checkmark \\ \hline
\cite{gomes_2025}       & \checkmark &            & \checkmark & \checkmark \\ \hline
\cite{han_2024}         & \checkmark & \checkmark &            & \checkmark \\ \hline
\cite{deng_2024}        &            & \checkmark &            & \checkmark \\ \hline
\cite{araujo_24}        &            & \checkmark &            & \checkmark \\ \hline
\cite{zanatta_23}       &            &            &            & \checkmark \\ \hline
\cite{hartikainen_2024} & \checkmark &            &            &            \\ \hline
\cite{wei_2025}         &            &            &            & \checkmark \\ \hline
\end{tabular}
\end{table}

Regarding the table \ref{tab:related_works}, we compiled all studies cited in this study related to observability research, highlighting their research gaps, limitations, and potential future work to develop a research roadmap.

\begin{table*}[!ht]
\centering
\caption{Analysis of Related Works and Primary Studies Cited}
\label{tab:related_works}
\footnotesize
% As larguras somam ~0.91\textwidth, deixando espaço para o padding (tabcolsep)
\begin{tabular}{p{0.15\textwidth} p{0.38\textwidth} p{0.38\textwidth}}
\toprule
\textbf{Ref.} & \textbf{Research Gaps and Limitations found} & \textbf{Future Works} \\ \midrule

Usman et al. \cite{usman_2022} &
Absence of formal reference architecture, operational complexity of fragmentation, telemetry limitations, and scalability and correlation challenges. & 
Intelligent Automation via AI/ML, Security integration, Enhancing Developer Experience, and Unified Observability Platforms. \\ \addlinespace

Araujo et al. \cite{araujo_24} &
Lack of holistic reference architectures, resource constraints (overhead) in fog nodes, and high heterogeneity of edge devices. & 
Lightweight AI/ML for local anomaly detection (TinyML), standardized interoperability protocols, and adaptive observability frameworks. \\ \addlinespace

Bruno et al. \cite{zanatta_23} &
Fragmented focus on metrics over logs/traces in 5G research, high memory/CPU overhead of standard CNOT stacks, and lack of standardized log structures. & 
Integration of distributed tracing, research on geographically distributed B5G setups, scalability analysis of COTS ML, and observability in Open RAN components. \\ \addlinespace

Kosińska et al. \cite{kosinska_2023} &
Lack of unified terminology, insufficient practical automation via AI/ML (AIOps gap), and poor adaptation of tools for edge computing environments. & 
Development of observability maturity models, deeper integration with eBPF for low-overhead visibility, and alignment with FinOps and GreenOps. \\ \addlinespace

Faseeha et al. \cite{faseeha_2025} &
Conceptual confusion between monitoring and observability pillars, lack of practical multi-pillar approaches, and difficulties in tracing failures in highly dynamic environments. & 
Evolution of frameworks for modern deployment paradigms (service mesh, event-driven), automation of data analysis, and optimization of observability overhead. \\ \addlinespace

Han et al. \cite{han_2024} &
Inconsistent problem definitions for RCA and high cloud-native heterogeneity across traces, metrics, and logs. & 
Improving multi-modal data correlation at scale and enhancing Graph Attention Networks for dynamic topologies. \\ \addlinespace

Gomes et al. \cite{gomes_2025} &
Difficulty in correlating multi-modal data (metrics, logs, traces), inadequacy of traditional monitoring for sophisticated CNAs, and complex root cause analysis (RCA) in distributed environments. & 
Advanced ML-based resource management, automated diagnostic mechanisms integrated with auto-scaling, and research on low-overhead real-time data collection. \\ \addlinespace

Wei et al. \cite{wei_2025} &
Lack of a systematic bridge between desirable system qualities and the concrete cloud-native practices and tools used to achieve them, and the initial validation scope is limited (expert-based, not covering all contexts). &
Provide mechanisms/interfaces to extend and continuously evolve the map (keeping practices/tools up to date) and perform broader validations beyond the initial expert review. \\ \addlinespace

Deng et al. \cite{deng_2024} &
Missing clear research roadmap for cloud-native lifecycles, high cognitive load for developers (DevEx), and insufficient real-time security observability. & 
Integration of SecOps with performance telemetry, improvement of Developer Experience (DevEx), and observability optimization for serverless and green computing. \\ \bottomrule
\end{tabular}
\end{table*}

To address the identified architectural gaps, we perceived in literature the possibility to introduce a formal representation of Observability ($O$) as a multidimensional set of telemetry data $O = \{M, L, T, E\}$, where $M, L, T, \text{and } E$ represent Metrics, Logs, Traces, and Events, respectively. Building on established maturity frameworks in the academic and industrial literature—such as CMMI, MPS.BR, ITIL, COBIT, and other models that outline progressive capability stages—we can devise an Observability Maturity Model (OMM) organized into five evolutionary levels (Fig. \ref{fig:obs-arch}): (1) Non-Existent, where no telemetry is captured; (2) Initial/Ad-Hoc, where data is collected but remains siloed in local storage; (3) Managed, characterized by centralized collection and aggregation for global system analysis; (4) Optimized, where visual analytics and dashboards provide actionable insights into system health; and (5) Predictive, the highest stage where AI-driven agents and machine learning models analyze the telemetry set $O$ to prevent failures before they impact the environment. Rather than introducing a novel structure, this model synthesizes and adapts patterns recurrently identified in prior studies on software process maturity, monitoring practices, and systems engineering.

\begin{figure}[ht]
    \centering
    \includegraphics[width=1.0\linewidth]{images/observability-architecture.png}
    \caption{Observability Maturity Model}
    \label{fig:obs-arch}
\end{figure}

To address the identified research gaps, Figure \ref{fig:research-road} presents a strategic research roadmap. We argue that to overcome the lack of formal standards and the high operational complexity, the next five years of research in Kubernetes observability must pivot from data collection toward three specific advancement pillars:

\begin{figure}[ht]
    \centering
    \includegraphics[width=1.0\linewidth]{images/research-roadpmap.png}
    \caption{Research Roadmap}
    \label{fig:research-road}
\end{figure}

\begin{enumerate}
    \item \textbf{Semantic Interoperability and Low-Overhead Instrumentation:} Transitioning from basic collection to deep, non-intrusive visibility using eBPF, ensuring that telemetry data from the kernel level to the application level is semantically correlated through OpenTelemetry standards.
    \item \textbf{AIOps for Multimodal Correlation:} Developing machine learning models capable of processing the ``Observability Tuple'' $O=\{M,L,T,E\}$ to perform automated root cause analysis (RCA) in highly dynamic and ephemeral environments, reducing the cognitive load on engineering teams.
    \item \textbf{Autonomic and Sustainable Observability (GreenOps):} Evolving toward ``Predictive'' maturity (Level 5), where observability data drives self-healing mechanisms and optimizes energy consumption in Kubernetes clusters, aligning performance with sustainability goals.
\end{enumerate}

\section{Conclusion}
\label{sec:conclusion}

This Rapid Review aimed to synthesize the state of the art of open-source architectures and tools for observability in cloud-native applications based on Kubernetes. Our findings confirm that the observability architecture is based on the pillars of metrics, logs, and traces, with events currently considered an additional pillar. Furthermore, we identified a gap in the literature regarding formal reference architectures, but we infer that a \textit{de facto} pattern pervades the implementation of these pillars.

As a Rapid Review, this study may not be as exhaustive as a full Systematic Literature Review (SLR). Furthermore, our deliberate focus on open-source solutions and English-language articles excludes insights that could be gained from grey literature, proprietary solutions, or non-English academic contributions. However, this focused approach aligns with our objective of providing timely, actionable guidance to practitioners seeking cost-effective observability solutions.

Regarding the specific objectives of this research, the study successfully drafted the two main contributions: 1) the introduction of a novel maturity observability model, and 2) a research roadmap to develop a robust open-source architecture for observability in a cloud-native environment using K8s.

By focusing on these specific challenges, the community can move beyond fragmented ``\textit{de facto}'' patterns toward a unified, intelligent, and architecture-driven observability ecosystem.

This review thus serves as a foundational guide for practitioners navigating the complex landscape of cloud-native observability and a roadmap for researchers aiming to bridge the identified gaps.\\

\textbf{Use of AI tools} \\
Large Language Models (LLMs) and AI-assisted tools, especially Gemini and Grammarly, were used solely for minor language editing and text polishing under the direct supervision of the authors. Images used in this paper were generated by Gemini and Napkin AI. AI generated none of the research design, interpretation, or scientific reasoning. The authors take full responsibility for the content and conclusions of this manuscript, and all conceptual and analytical work was conducted entirely by human authors.

\section{Declarations}
\label{sec:declarations}

\textbf{Funding}: Not applicable. \\
\textbf{Conflicts of Interest}: Not applicable.\\
\textbf{Ethical approval and Consent to Participate}: Not applicable, since this study does not involve human participants or animals. \\
\textbf{Consent for Publication}: Not applicable. \\
\textbf{Data Availability}: Not applicable since no external data were generated in this phase of this study.\\
\textbf{Materials Availability}: Not applicable.\\
\textbf{Author Contributions}: Author 1 conducted the study design, search, screening, and synthesis. Author 2 guided the methodology and performed a critical review of the text. Author 3 supervised the study and contributed to consolidating the discussion and conclusion sections. All authors read, reviewed, and approved the final version of the manuscript.

%% The next two lines define the bibliography style to be used, and
%% the bibliography file.
\bibliographystyle{ACM-Reference-Format}
\bibliography{sample-base}

%% If your work has appendices, this is the place to put them
%\appendix

\end{document}

% ============================================================
% CBSoft ACM-like Template – Author Checklist
% ============================================================
%
% 1. Remove CCS Concepts
%    - Remove the CCSXML block and \ccsdesc commands.
%
% 2. Remove ACM footer
%    Use the command \renewcommand\footnotetextcopyrightpermission[1]{}
%
% 3. Remove ACM Reference Format
%    Use the commands \setcopyright{none} and 
%    \settopmatter{printacmref=false}
%
% 4. Short title must not exceed column width
%    Use the command 
%    \renewcommand{\shorttitle}{The Name of the Title Is Hope}
%
% 5. Include authors' surnames in the page headers (even pages)
%    - The author names in the page header (at the top of the page)
%      must not exceed the column width
%    - This list must contain only the authors' surnames
%    - Suggestion: use "Surname et al." for more than two authors with the
%      command \renewcommand{\shortauthors}{Surname1stAuthor et al.}
%
% 6. Author format below title:
%    \author{FirstName Surname}
%    \affiliation{
%      \institution{Department, Institution}
%      \city{City}
%      \country{Country}}
%    \email{email@email.com}
%
% 7. Verify all author emails are included
%    - Check font consistency
%
% 8. All numbered section and subsection titles must follow the same
%    capitalization style: the first letter of each word in the title 
%    must be capitalized (Title Case).
%
% 9. Artifact Availability (check CFP)
%    - Not numbered -- use \section*{Artifact Availability}
%    - Placed after Conclusion and before Acknowledgments
%    - In Portuguese, use \section*{Disponibilidade de Artefatos}
%
% 10. Abstract
%     - Single paragraph
%     - Papers written in Portuguese must use "Resumo" instead of "Abstract"
%       with the command \renewcommand\abstractname{Resumo}
%
% 11. Keywords
%     - Comma-separated, first letter capitalized
%     - Papers written in Portuguese must use "Palavras-chave" instead of
%       "Keywords" with the command 
%       \renewcommand\keywordsname{Palavras-chave}
%
% 12. References
%     - Not numbered -- already enforced by the template
%     - Papers written in Portuguese must use "Referências" instead of
%       "References" with the command \renewcommand\refname{Referências}
%
% 13. Acknowledgements
%     - Optional section
%     - Not numbered -- use \section*{Acknowledgements}
%     - In Portuguese, use \section*{Agradecimentos}
%
% 14. Remove received/revised/accepted dates
%
% 15. Figure captions must appear below figures
%
% 16. Verify the symposium name in the page header
%    - Use English, even if the paper is written in Portuguese
%    - Select the correct event according to the track
%    - Use Arabic ordinal numbers, not Roman
%
%    SBES (Research, Education, IIER, Tools, Industry)
%    \acmConference[SBES 2026]{40th Brazilian Symposium on Software Engineering}{September 8--11, 2026}{São Paulo, SP, Brazil}
%
%    SBLP
%    \acmConference[SBLP 2026]{30th Brazilian Symposium on Programming Languages}{September 8--11, 2026}{São Paulo, SP, Brazil}
%
%    SBCARS
%    \acmConference[SBCARS 2026]{20th Brazilian Symposium on Software Components, Architectures, and Reuse}{September 8--11, 2026}{São Paulo, SP, Brazil}
%
%    SAST
%    \acmConference[SAST 2026]{11st Brazilian Symposium on Systematic and Automated Software Testing}{September 8--11, 2026}{São Paulo, SP, Brazil}
%
%    Workshops
%    \acmConference[Workshop Acronym 2026]{Workshop Full Name}{Workshop Date}{São Paulo, SP, Brazil}
%
%
% 17. Remove colored author names, citations, and links
%
% 18. Camera-ready only:
%     Do not use the "review" or "anonymous" modes in the 
%     \documentclass command
%
% 19. SBES Tools Track papers must include a demo video link at the end of
%     the abstract
%     Example:
%     \begin{abstract}
%     <abstract text> \newline
%     \textbf{Demo video:} <video link>
%     \end{abstract}
%
% 20. SBES Industry Track papers must not include author biographies
%
% ============================================================

\endinput


